\documentclass[12pt,a4paper]{article}  % Déclare la classe du document.

\newcommand{\chaptermark}[2]{\og #2 \fg{} (#1)} % ligne uniquement pour éviter une erreur en cas d'article
\usepackage{mathrsfs}
\include{../1ere-preambule}

\usepackage{epic}   % extension et simplification de  l'environnement picture
\usepackage{graphicx}    % permet d'inclure des graphique externe dans un document
\usepackage{arydshln}   %permet d'avoir des lignes en pointillés

\newcommand{\lignes}[1]{
\setlength{\unitlength}{1mm}
\vskip 0,4cm\noindent
\multido{\i=0+1}{#1}{%
     \noindent
     \begin{picture}(150,5)
     \linethickness{0,005mm}
     \put(-5,0){\dashline{1}(175,5)(5,5)}
     \end{picture}\vskip 0,001cm
     }
\vskip 0cm
}

\rhead{DS n$^{\circ}$06 - Sujet A - 15/05/2022}

\begin{document}
\graphicspath{{../images/}}

\begin{center}
\fbox{
\begin{minipage}[]{15cm}
\begin{center}
\LARGE{\rule{0.5cm}{0pt}\rule[-0.25cm]{0pt}{0.9cm}
DS n$^{\circ}$06 - Sujet A - Fonction exponentielle\\
Application du produit scalaire
}
\end{center}
\end{minipage}
}

\end{center}


\begin{center}\textbf{1h - Calculatrice autorisée - Barème indicatif}\end{center}

\begin{center}\textbf{Les élèves avec un tier-temps ne traitent pas les questions avec  le symbole \ding{46}}\end{center}

\begin{center}
%On veillera à ENCADRER ses résultats et soigner sa copie.\\
Les résultats doivent être justifiés par des calculs (au moins 1 étape intermédiaire).
\end{center}

\begin{center}\textbf{Les élèves avec un tier-temps ne traitent pas les questions avec  le symbole \ding{46}}\end{center}

\exerciceDS{ - Au golf - ( 2 points )}
\begin{minipage}[c]{.6\linewidth}
Sur un green de golf, trois positions de la balle sont repérées en $\mathrm{M},$ N et P sur le schéma ci-contre. Le triangle MNP est tel que MN=8, $\mathrm{MP}=5$ et $\widehat{M}=53^{\circ}$.\\
Calculer la valeur exacte de NP, puis son arrondi à 0,1 près.
\end{minipage}
\hfill
\begin{minipage}[c]{.35\linewidth}
\begin{center}
   \includegraphics[scale=0.8]{1.05.DS-IMG-01} 
\end{center}

\end{minipage}

\exerciceDS{ - Angles dans un triangle - ( 3 points )}
Dans un repère orthonormé du plan, on considère les points $\mathrm{A}(-2 ;3), \mathrm{B}(1 ; 4)$ et $\mathrm{C}(4 ; 1)$.\\
Déterminer les angles du triangle $A B C$ à 0,1 degré près.\\


\exerciceDS{ - Médiatrice dans un triangle - ( 8 points )}
On considère le triangle $\mathrm{ABC}$ tel que $\mathrm{A}(-2 ;3), \mathrm{B}(1 ; 4)$ et $\mathrm{C}(4 ; 1)$.
\begin{enumMet}
    \item Montrer que les coordonnées du milieu M du segment $[BC]$ sont $M(2,5;2,5)$
    \item Déterminer une équation de la médiatrice  du segment $[BC]$ dans le triangle $\mathrm{ABC}$ (la médiatrice est la perpendiculaire à $[BC]$ passant par M).
    \item \ding{46} Montrer que les coordonnées du milieu N du segment $[AC]$ sont $N(1;2)$
    \item Déterminer une équation de la médiatrice  du segment $[AC]$ dans le triangle $\mathrm{ABC}$.
    \item Déterminer les coordonnées du point d'intersection $\mathrm{O}$ des deux médiatrices.
    \item \ding{46} Montrer que les coordonnées du milieu P du segment $[AB]$ sont $P(-0,5;3,5)$
    \item Calculer le produit scalaire $\vv{\mathrm{AB}} \cdot \vv{\mathrm{OP}}$. Quel résultat connu retrouve-t-on ?
\end{enumMet}


\exerciceDS{ - Équations cartésiennes - ( 3 points )}
Dans chaque cas, déterminer une équation cartésienne de la droite $d$ passant par le point $\mathrm{A}$ et perpendiculaire à la droite (BC).
\begin{enumMet}
    \item \ding{46} $\mathrm{A}(5 ; 4)$ et $\Coord{u}{-2}{4}$ un vecteur directeur de $(\mathrm{BC})$.
    \item $\mathrm{A}(0 ; 2), \mathrm{B}(3 ;-1)$ et $\mathrm{C}(-4 ; 2)$.
    \item $\mathrm{A}(-2 ; 3)$ et la droite $(\mathrm{BC})$ est parallèle à l'axe des abscisses.
\end{enumMet}

\exerciceDS{ - simplification - ( 4 points )}
$x$ est un réel quelconque. Simplifier au maximum les expressions suivantes.
\begin{multicols}{2}
	\begin{enumMet}
	\item $a(x)=\mathrm{e}^{2 x} \times\left(\mathrm{e}^x\right)^2 \times \mathrm{e}^{-3 x}$
	\item $b(x)=\dfrac{\mathrm{e}^{x^2}}{\mathrm{e}^x}$
	\item $c(x)=\dfrac{\mathrm{e}^{x-1} \times \mathrm{e}^{4 x}}{\mathrm{e}^x}$
	\item \ding{46} $d(x)=\dfrac{\mathrm{e}^{-2 x}}{\mathrm{e}^{-3 x} \times \mathrm{e}^{x+1}}$
\end{enumMet}
\end{multicols}

\exerciceDS{ - dérivation - ( 4 points )}
Calculer la dérivée de la fonction $f$ définie sur $\mathbb{R}$.
\begin{multicols}{2}
	\begin{enumMet}
		\item $f(x)=4 \mathrm{e}^x-8 x^2+2$
		\item \ding{46} $f(x)=2x-3 \mathrm{e}^x+1$
		\item $f(x)=2\mathrm{e}^x+\mathrm{e}^3$
		\item $f(x)=2x \mathrm{e}^x$
	\end{enumMet}
\end{multicols}

\exerciceDS{ - Résolution d'équations et d'inéquations - ( 8 points )}
Résoudre dans \R.
\begin{multicols}{2}
	\begin{enumMet}
		\item $\mathrm{e}^{2x-3}=\mathrm{e}$
		\item $\mathrm{e}^{-2x^2}=\dfrac{1}{\mathrm{e}^{-8}}$
		\item \ding{46} $\mathrm{e}^{2x^2-6x}=1$
		\item $5+\mathrm{e}^x=1$
		\item $\mathrm{e}^{x+1}<1$
		\item $\mathrm{e}^{-2 x+5} \geqslant 0$
		\item $-3 \mathrm{e}^{x^2-4}>4$
		\item $\mathrm{e}^{x+4} \leqslant \dfrac{1}{\mathrm{e}^{3 x}}$
	\end{enumMet}
\end{multicols}



\end{document}